\section{The four-point residue: fusion, area, and an explicit density}
\label{sec:screenedfusion}

The one-screening calculation can be completed at the level of
four-point residues. Two intermediate Virasoro channels contribute,
their change of basis is an explicit hypergeometric connection
matrix, and the invariant Hermitian pairing is exactly the flat
area. This gives both a density and an evaluation of its total
volume for arbitrary four deficit weights. At balanced weights,
the two ordinary blocks coalesce and must be replaced by a
regularized logarithmic basis.

\subsection{The two spectral pinches and their coefficients}

Put the fields with charges $\alpha_1,\alpha_2,\alpha_3,\alpha_4$
at $0,1,z,\infty$, respectively. Throughout this section
\begin{equation}
 0<b<1,\qquad Q=b+b^{-1},\qquad
 \alpha_i=\frac{\mu_i}{2b},\qquad
 0<\mu_i<1,\qquad \sum_i\mu_i=2.
 \label{eq:screeningparameters}
\end{equation}
To take a transverse residue, vary the charges near this locus
and write
\begin{equation}
 \varepsilon=\sum_i\alpha_i-Q+b,\qquad
 a_s=\alpha_1+\alpha_3.
 \label{eq:screeningtransverse}
\end{equation}
The $s$-channel pairs the marks $0,z$; its spectral integrand
contains
\[
 C(\alpha_1,\alpha_3,Q-\alpha)
 C(\alpha,\alpha_2,\alpha_4)
 \mathcal F_\alpha(z)\mathcal F_\alpha(\bar z).
\]
For real charges the last product is the absolute square.
For continuation in complex charges, both chiral factors are
continued analytically, rather than conjugating the variable
charges. We use the DOZZ/KRV normalization fixed in
Section~\ref{sec:quantummeasure}.

The bootstrap identity is established in an initial convergence
domain by~\cite[Theorem~1.1]{GKRVBootstrap}. Its continuation
requires the discrete terms created by crossing poles, as the
paragraph following that theorem also explains. The relevant
contour mechanism and pole families are described in
\cite[Section~7.2, eqs.~(87)--(90)]{TeschnerRevisited}.

\paragraph{Analytic input for the spectral derivation.}
The infinite integral requires a separate continuation statement.
In the present charge convention and ordering, the initial real
bootstrap domain has $\alpha_i<Q/2$,
$\alpha_1+\alpha_3>Q/2$, and $\alpha_2+\alpha_4>Q/2$.
The one-screening locus $\sum_i\alpha_i=Q-b$ lies outside that
domain. Meromorphic dependence of each integrand factor alone
does not justify continuing the integral. The spectral derivation
below uses the following sufficient input, at fixed $0<b<1$.

Continue from the initial domain along a specified path of complex
charges, retaining every residue of a pole crossed by the contour.
Locally, the deformed contours are required to coincide with
$Q/2+i\RR$ outside a fixed compact set. Near a generic target point,
separate the finitely many pinch contributions, including their
reflected partners, and write the remaining term using a locally
fixed contour $\Gamma_{\rm reg}$ avoiding all poles. Its integrand
$\mathcal I_{\rm reg}(\alpha;\varepsilon,\eta,z)$ must be holomorphic
in the charge coordinates through $\varepsilon=0$; here $\eta$
denotes the remaining charge coordinates. Any additional discrete
terms left in this remainder must also be holomorphic there.
Require normal absolute convergence: on every compact set $K$ of
these complex charge parameters and of noncolliding $z$ in the
chosen branch chart, there is $G_K\in L^1(\RR)$ such that on the
common contour ends
\begin{equation}
 \sup_{(\varepsilon,\eta,z)\in K}
 \left|\mathcal I_{\rm reg}(Q/2+iP;\varepsilon,\eta,z)\right|
 \le G_K(P).
 \label{eq:screeningtailinput}
\end{equation}
The same normal-convergence requirement is imposed on compact
continuation neighborhoods along the path, after their crossed-pole
terms are separated. If another contour prescription moves the ends,
the connecting pieces at infinity must additionally have integrals
tending to zero uniformly on those parameter sets.

Under these conditions, the truncated regular integrals are
holomorphic in $\varepsilon$ and converge locally uniformly as the
spectral cutoff tends to infinity. Their limit is therefore
holomorphic through $\varepsilon=0$, and Cauchy's residue integral
in that parameter commutes with removal of the cutoff. The regular
term contributes no residue; only the finite pinch terms do.

This is an additional analytic input for the contour proof. No
estimate~\eqref{eq:screeningtailinput} is proved here, and the cited
bootstrap theorem in its initial domain is not being used as such
an estimate at the one-screening locus. The residue identity below
also has an independent first-moment and period derivation. The
curvature identity and finite fusion formulas use that derivation.
Balanced limits are taken in the finite period expressions; no
uniform spectral estimate through further pole coalescences or as
$b\to0$ is asserted.

\paragraph{The local pinch calculation.}
Two opposite pole families are
\begin{align}
 \alpha&=a_s+mb+n/b,\notag\\
 \alpha&=Q-\alpha_2-\alpha_4-m'b-n'/b
        =a_s+b-\varepsilon-m'b-n'/b,
 \label{eq:screeningpolefamilies}
\end{align}
with nonnegative integer indices. At $\varepsilon=0$ they collide
only if
\[
 (m+m'-1)b^2+(n+n')=0.
\]
Since $0<b^2<1$, this forces $m+m'=1$ and $n=n'=0$.
Thus the two intermediate charges are
\begin{equation}
 \boxed{\quad \alpha_s^{(0)}=a_s,\qquad
                 \alpha_s^{(1)}=a_s+b.\quad}
 \label{eq:twoscreeningchannels}
\end{equation}
Their reflected partners $Q-\alpha_s^{(j)}$ are included in the
usual unfolded contour. We first work at generic weights and
then take the regularized limits at coincident poles.

The coefficients require only elementary special-function
identities. Define
\begin{equation}
 \mathfrak g(x)=\frac{\Gamma(x)}{\Gamma(1-x)}.
 \label{eq:gammaratio}
\end{equation}
For a three-point function with $\Sigma=\alpha_1+\alpha_2+\alpha_3$,
the zero- and one-screening residues, with respect to the
displayed sum coordinate, are
\begin{align}
 \operatorname*{Res}_{\Sigma=Q}C&=1,\notag\\
 \operatorname*{Res}_{\Sigma=Q-b}C
  &=-\frac{\pi\Lambda_L}
       {\prod_{j=1}^3\mathfrak g(2b\alpha_j)}.
 \label{eq:threepointscreeningresidues}
\end{align}
For completeness these follow directly from DOZZ. Its denominator
is $\Upsilon_b(\Sigma-Q)\prod_j\Upsilon_b(\Sigma-2\alpha_j)$;
use
\[
 \Upsilon_b(Q-x)=\Upsilon_b(x),\qquad
 \Upsilon_b(x+b)=\mathfrak g(bx)b^{1-2bx}\Upsilon_b(x).
\]
At $\Sigma=Q$ reflection cancels the other factors. At
$\Sigma=Q-b=1/b$, the product of shifted ratios is
$b/\prod_j\mathfrak g(2b\alpha_j)$. Combining it with the
DOZZ prefactor and
$\mathfrak g(b^2)\mathfrak g(-b^2)=-b^{-4}$ proves the
second identity. These are the conventions of
\cite[eqs.~(9), (11), (13)]{TeschnerRevisited}.

Set
\begin{equation}
 \delta_0=1-\mu_1-\mu_3,\qquad
 \delta_1=1-\mu_2-\mu_3,\qquad
 \delta_\infty=1-\mu_3-\mu_4,
 \label{eq:screeningexponentdifferences}
\end{equation}
and, initially assuming $\delta_0\ne0$, define
\begin{equation}
 h_0=\frac{\pi\mathfrak g(\delta_0)}
              {\mathfrak g(\mu_2)\mathfrak g(\mu_4)},
 \qquad
 h_1=\frac{\pi\mathfrak g(-\delta_0)}
              {\mathfrak g(\mu_1)\mathfrak g(\mu_3)}.
 \label{eq:screeningHermitiancoefficients}
\end{equation}
At $\alpha=a_s$, the left three-point function has its
zero-screening pole and the right one its one-screening pole,
with coefficient $-\Lambda_Lh_0$. At $\alpha=a_s+b$ their
roles reverse, with coefficient $-\Lambda_Lh_1$.
In the local coordinate $x=\alpha-\alpha_s^{(j)}$ the singular
product is
\[
 \frac{\Lambda_Lh_j}{x(x+\varepsilon)}.
\]
The contour separating $x=-\varepsilon$ and $x=0$ gives the
coefficient $-\Lambda_Lh_j/\varepsilon$. In the unfolded
$1/(4\pi i)$ spectral normalization the reflection partner
supplies the other half of the local $1/(2\pi i)$ residue.
Under the continuation and tail-control input above, the remaining
contour integral is holomorphic in $\varepsilon$ and contributes
no residue. The contour calculation therefore gives
\begin{equation}
 \boxed{\quad
 \operatorname*{Res}_{\varepsilon=0}C_4(z)
  =-\Lambda_L\left(
      h_0|\mathcal F_0(z)|^2+h_1|\mathcal F_1(z)|^2\right).
 \quad}
 \label{eq:spectralonescreeningresidue}
\end{equation}
The identity itself also follows without continuing the infinite
spectral integral: combine the first-moment residue
\eqref{eq:onescreeningresidue}, in the transverse coordinate
$\varepsilon=b(s+1)$, with the independently proved period pairing
\eqref{eq:screeningareapairing} and the block factorization
\eqref{eq:screeningVirasoroblocks}. This fixes the normalization
and sign independently of the conditional contour derivation.
Conventions for the overall normalization of correlations differ
in the literature; the bootstrap normalization is calibrated by
\eqref{eq:threepointscreeningresidues}.

\subsection{Why the screening integrals are the Virasoro blocks}

Let $F(a,b;c;z)={}_2F_1(a,b;c;z)$ and take branches positive
for $0<z<1$, continued first to
$\CC\setminus(({-\infty},0]\cup[1,\infty))$.
The normalized $s$-channel periods are
\begin{align}
 f_0(z)&=F(1-\mu_4,\mu_3;\mu_1+\mu_3;z),\notag\\
 f_1(z)&=z^{\delta_0}
        F(\mu_2,1-\mu_1;1+\delta_0;z).
 \label{eq:screeningsbasis}
\end{align}
The full normalized chiral blocks are
\begin{equation}
 \mathcal F_j(z)=E_b(z)f_j(z),\qquad
 E_b(z)=z^{-\mu_1\mu_3/(2b^2)}
                  (1-z)^{-\mu_2\mu_3/(2b^2)}.
 \label{eq:screeningVirasoroblocks}
\end{equation}
Here the conformal weight is $\Delta(\alpha)=\alpha(Q-\alpha)$
and the central charge is $c_L=1+6Q^2$.

To verify the block identification, the screening field has
weight $\Delta(b)=b(Q-b)=1$. Integration over a closed twisted
cycle therefore commutes with the Virasoro action: its Ward
commutator is a total derivative. A screening outside a circle
separating $0,z$ from $1,\infty$ gives intermediate charge
$a_s$, while a screening inside gives $a_s+b$.
The corresponding Euler arc integrals are
\begin{align}
 S_0&=\int_1^\infty
   w^{-\mu_1}(w-1)^{-\mu_2}(w-z)^{-\mu_3}\dd w
       =B(1-\mu_2,1-\mu_4)f_0(z),\notag\\
 S_1&=\int_0^z
   w^{-\mu_1}(1-w)^{-\mu_2}(z-w)^{-\mu_3}\dd w
       =B(1-\mu_1,1-\mu_3)f_1(z).
 \label{eq:screeningEulerarcs}
\end{align}
The positive arc branches are phase-normalized versions of
the twisted periods; closed Pochhammer contours justify the
Ward-identity argument, with the same normalized functions.
Dividing out the beta constants sets the leading coefficient
to one. The leading exponents agree because
\[
 \Delta(a_s)-\Delta(\alpha_1)-\Delta(\alpha_3)
   =-2\alpha_1\alpha_3,
 \qquad
 \Delta(a_s+b)-\Delta(a_s)=\delta_0.
\]
The screened intertwiner construction and uniqueness of the
normalized Ward solution identify the functions at every
level; see~\cite[Section~18.1, eq.~(224)]{TeschnerRevisited}.
Here $0<a_s,a_s+b<Q$, so these internal weights are positive
and avoid the degenerate Verma weights at $c_L>25$.
No external degenerate field is assumed.

The two blocks solve a second-order hypergeometric equation
after removal of the common factor $E_b$. This reduction is a
consequence of the specified one-screening charge locus, not
a general second-order equation for arbitrary four-point
Liouville blocks.

\subsection{The finite fusion matrix}

For $\delta_1\ne0$, the normalized $t$-channel periods, pairing
$z,1$, are
\begin{align}
 g_0(z)&=F(1-\mu_4,\mu_3;\mu_2+\mu_3;1-z),\notag\\
 g_1(z)&=(1-z)^{\delta_1}
        F(1-\mu_2,\mu_1;1+\delta_1;1-z).
 \label{eq:screeningtbasis}
\end{align}
Their internal charges are $\alpha_2+\alpha_3$ and
$\alpha_2+\alpha_3+b$, and their full chiral blocks are
$E_b g_j$. Thus the same matrix relates periods and normalized
screened Virasoro blocks:
\begin{equation}
 \begin{pmatrix}f_0\\f_1\end{pmatrix}
  =\mathsf M\begin{pmatrix}g_0\\g_1\end{pmatrix},
 \label{eq:screeningfusionrelation}
\end{equation}
where
\begin{equation}
 \boxed{\quad
 \mathsf M=\begin{pmatrix}
 \dfrac{\Gamma(1-\delta_0)\Gamma(\delta_1)}
       {\Gamma(\mu_1)\Gamma(1-\mu_2)}
 &\dfrac{\Gamma(1-\delta_0)\Gamma(-\delta_1)}
       {\Gamma(\mu_3)\Gamma(1-\mu_4)}\\[10pt]
 \dfrac{\Gamma(1+\delta_0)\Gamma(\delta_1)}
       {\Gamma(1-\mu_3)\Gamma(\mu_4)}
 &\dfrac{\Gamma(1+\delta_0)\Gamma(-\delta_1)}
       {\Gamma(1-\mu_1)\Gamma(\mu_2)}
 \end{pmatrix}.
 \quad}
 \label{eq:screeningfusionmatrix}
\end{equation}
This is the Gauss connection formula, obtained by substituting
$a=1-\mu_4$, $b=\mu_3$, $c=\mu_1+\mu_3$ in
\cite[eqs.~15.10.21--22]{DLMF}. The Wronskians give
\begin{equation}
 \det\mathsf M=-\frac{\delta_0}{\delta_1}.
 \label{eq:fusiondeterminant}
\end{equation}
The matrix is independent of the Liouville parameter $b$ when
the deficits $\mu_i$ are fixed. All $b$-dependence in these
blocks is then in the common factor $E_b$ and in the overall
correlation normalization.

There is a simpler expression in actual arc coordinates.
For real $0<z<1$, put
\begin{align*}
 T_0&=\int_{-\infty}^0
  (-w)^{-\mu_1}(1-w)^{-\mu_2}(z-w)^{-\mu_3}\dd w
       =B(1-\mu_1,1-\mu_4)g_0,\\
 T_1&=\int_z^1
  w^{-\mu_1}(1-w)^{-\mu_2}(w-z)^{-\mu_3}\dd w
       =B(1-\mu_2,1-\mu_3)g_1.
\end{align*}
Gamma reflection reduces~\eqref{eq:screeningfusionmatrix} to
\begin{equation}
 \boxed{\quad
 \begin{pmatrix}S_0\\S_1\end{pmatrix}
 =\frac1{\sin(\pi\delta_1)}
 \begin{pmatrix}
 \sin(\pi\mu_1)&-\sin(\pi\mu_3)\\
 \sin(\pi\mu_4)&-\sin(\pi\mu_2)
 \end{pmatrix}
 \begin{pmatrix}T_0\\T_1\end{pmatrix}.
 \quad}
 \label{eq:screeningsinefusion}
\end{equation}
These arc pairs cease to be bases at the corresponding
balanced weights. The divergent coefficients in this formula
cannot be specialized there without first changing basis.

\subsection{The Hermitian pairing and its direct verification}

We now prove the area pairing directly from periods:
\begin{equation}
 \boxed{\quad
 \mathcal A_\mu(z)=h_0|f_0(z)|^2+h_1|f_1(z)|^2.
 \quad}
 \label{eq:screeningareapairing}
\end{equation}
Together with the first-moment residue, this supplies an
independent derivation of~\eqref{eq:spectralonescreeningresidue}.
The complex beta integral is
\begin{equation}
 \int_{\CC}|w|^{-2p}|1-w|^{-2q}\dd^2w
 =\pi\frac{\Gamma(1-p)\Gamma(1-q)\Gamma(p+q-1)}
             {\Gamma(p)\Gamma(q)\Gamma(2-p-q)},
 \label{eq:complexbetaintegral}
\end{equation}
for $0<p,q<1$ and $p+q>1$. Introduce positive Laplace
parameters for the two powers, perform the Gaussian integral
in $w$, and then use their sum and ratio as variables. The
remaining integrals are an ordinary Gamma integral and
$B(1-q,1-p)$, proving the formula.

The area is a constant Hermitian form in the period basis.
For $\delta_0\ne0$, local monodromy has distinct eigenvalues
$1,e^{2\pi i\delta_0}$, so the form is diagonal in
$(f_0,f_1)$. If $\delta_0>0$, sending $z\to0$ in the area
integral and using~\eqref{eq:complexbetaintegral} gives $h_0$.
If $\delta_0<0$, rescaling $w=zu$ gives the leading
coefficient $h_1$. To determine the other coefficient, first
assume also $\delta_1\ne0$. Monodromy at $1$ makes the
Hermitian form diagonal in $(g_0,g_1)$, so the off-diagonal
entry after the real connection matrix vanishes:
\[
 h_0\mathsf M_{00}\mathsf M_{01}
       +h_1\mathsf M_{10}\mathsf M_{11}=0.
\]
Every matrix entry here is nonzero at these real nonresonant
weights. Hence either already determined coefficient fixes
the other; Gamma reflection gives precisely
\eqref{eq:screeningHermitiancoefficients}. Continuity at
$\delta_1=0$, with $\delta_0\ne0$ fixed, completes this
argument without crossing the singular $s$-channel basis.
For this continuity the area integrals are locally dominated:
on compact sets of admissible weights their singular exponents
stay below $2$, and their decay at infinity stays faster than
$|w|^{-2}$.

The signs of $h_0,h_1$ are opposite. The pairing has signature
$(1,1)$; the positive value of $\mathcal A_\mu$ on the period
line is compatible with that indefinite signature. It is not
a positive probability decomposition into two channels.
In arc coordinates it becomes
\begin{equation}
 \mathcal A_\mu=
 \frac{\sin(\pi\mu_2)\sin(\pi\mu_4)|S_0|^2
       -\sin(\pi\mu_1)\sin(\pi\mu_3)|S_1|^2}
      {\sin(\pi\delta_0)}.
 \label{eq:screeningarcarea}
\end{equation}
Define $H_s=\diag(h_0,h_1)$ and
\[
 H_t=\diag\left(
   \frac{\pi\mathfrak g(\delta_1)}
        {\mathfrak g(\mu_1)\mathfrak g(\mu_4)},
   \frac{\pi\mathfrak g(-\delta_1)}
        {\mathfrak g(\mu_2)\mathfrak g(\mu_3)}\right).
\]
The connection preserves the area exactly:
\begin{equation}
 \mathsf M^*H_s\mathsf M=H_t.
 \label{eq:fusionHermitianidentity}
\end{equation}
This follows either from the two area decompositions or
directly from Gamma recurrence and reflection. On the
specified real branch $\mathsf M^*=\mathsf M^T$.

\subsection{A closed formula for the four-point density and total volume}

The area expression gives an explicit density in the conformal
cross-ratio. In this section the period notation is
$\omega_{\mathrm{DM},\mu}=\omega_{\mathrm{CH},\mu}$ and
$V_{\mathrm{DM}}(\mu)=V_{\mathrm{CH}}(\mu)$, with the same
curvature-$-1$ normalization as~\eqref{eq:CHpotential}.
For a holomorphic column $\boldsymbol f=(f_0,f_1)^T$
and constant Hermitian matrix $H$, differentiation gives
\begin{equation}
 \partial_z\partial_{\bar z}\log(\boldsymbol f^*H\boldsymbol f)
 =\frac{\det H\,|f_0f_1'-f_0'f_1|^2}
        {(\boldsymbol f^*H\boldsymbol f)^2}.
 \label{eq:HermitianWronskian}
\end{equation}
The hypergeometric equation yields
\begin{equation}
 W(f_0,f_1)=\delta_0
      z^{-(\mu_1+\mu_3)}(1-z)^{-(\mu_2+\mu_3)},
 \qquad
 -h_0h_1\delta_0^2
       =\frac{\pi^2}{\prod_i\mathfrak g(\mu_i)}.
 \label{eq:screeningWronskian}
\end{equation}
Using $\dd z\wedge\dd\bar z=-2i\dd x\wedge\dd y$ and
the curvature-$-1$ convention therefore gives
\begin{equation}
 \boxed{\quad
 \omega_{\mathrm{DM},\mu}
 =\frac{4\pi^2}{\prod_{i=1}^4\mathfrak g(\mu_i)}
  \frac{|z|^{-2(\mu_1+\mu_3)}
        |1-z|^{-2(\mu_2+\mu_3)}}{\mathcal A_\mu(z)^2}
       \dd x\wedge\dd y,\qquad z=x+iy.
 \quad}
 \label{eq:general4DMdensity}
\end{equation}
The area in the denominator is explicitly the hypergeometric
combination~\eqref{eq:screeningareapairing}. Although that
diagonal basis is singular at balanced weights, the density
extends through them by its logarithmic limit.

Its integral can also be evaluated. At $z=0$, if $\delta_0>0$
then $\mathcal A_\mu\to h_0$, while if $\delta_0<0$ then
$\mathcal A_\mu\sim h_1|z|^{2\delta_0}$. In both cases
the metric density has the form
$\text{const.}\,|z|^{-2(1-|\delta_0|)}$, giving cone
angle $2\pi|\delta_0|$. At $\delta_0=0$ the area has
logarithmic growth and the end is a cusp. The same argument
at $1,\infty$ gives the other two angles
$2\pi|\delta_1|$ and $2\pi|\delta_\infty|$.
This agrees with the three-ended cone-sphere completion
described in~\cite[Section~5]{Kojima}.

Gauss--Bonnet, retaining the labels, gives the integral
\begin{equation}
 \boxed{\quad
 V_{\mathrm{DM}}(\mu)
 =2\pi\bigl(1-|\delta_0|-|\delta_1|-|\delta_\infty|\bigr)
 =4\pi\min_{1\le i\le4}\{\mu_i,1-\mu_i\}.
 \quad}
 \label{eq:general4DMvolume}
\end{equation}
For the last equality, order the weights as
$x_1\le x_2\le x_3\le x_4$ and put $r=1-x_4$.
The sum constraint forces $x_2,x_3\ge r$. The three absolute
bipartition differences have sum
\[
 \sum_{j=1}^3|r-x_j|
   =1-r-x_1+|r-x_1|=1-2\min(r,x_1).
\]
This proves the stated minimum formula. Equal weights give
$V_{\mathrm{DM}}=2\pi$; for example
$\mu=(1/4,1/2,1/2,3/4)$ gives $V_{\mathrm{DM}}=\pi$.

The integral here is of the smooth curvature on the open
moduli space, with its cone/cusp completion. If one instead
extends $\partial\bar\partial\log C_b^{\mathrm{uv}}$ as a
distribution to the compactified sphere, the explicit
pairwise powers have delta masses at collisions. Their
logarithms are pluriharmonic only off those collisions.

\subsection{Four equal weights: the logarithmic fusion calculation}

When all $\mu_i=1/2$, both $\delta_0$ and $\delta_1$ vanish.
The functions $f_0,f_1$ coincide. On the representation side,
$a_s+b=Q-a_s$, so the two intermediate charges have the
same Virasoro weight. The second solution is a regularized
logarithmic, or parameter-derivative, block; it is not a
second independent ordinary normalized highest-weight block
at the same weight.

Use the independent periods
\begin{equation}
 f(z)=F(1/2,1/2;1;z)=\frac2\pi K(z),\qquad
 g(z)=F(1/2,1/2;1;1-z)=\frac2\pi K(1-z).
 \label{eq:equalweightfg}
\end{equation}
Here $K$ uses the parameter convention of
\eqref{eq:ellipticperiodratio}. Put $L=\log16$ and define
\begin{equation}
 \ell_0=Lf-\pi g,\qquad \ell_1=Lg-\pi f.
 \label{eq:equalweightlogsolutions}
\end{equation}
The logarithmic expansion~\cite[Sections~15.10 and 19.12]{DLMF}
gives
\[
 \ell_0(z)=f(z)\log z+
   4\sum_{m\ge1}
    \left(\frac{(1/2)_m}{m!}\right)^2
       (H_{2m}-H_m)z^m,
\]
where $H_m$ is a harmonic number. Thus the holomorphic part
has zero constant term. The analogous statement for
$\ell_1$ uses $\log(1-z)$.
Solving~\eqref{eq:equalweightlogsolutions} gives the finite
connection matrix
\begin{equation}
 \boxed{\quad
 \begin{pmatrix}f\\\ell_0\end{pmatrix}
 =\frac1\pi
 \begin{pmatrix}L&-1\\L^2-\pi^2&-L\end{pmatrix}
 \begin{pmatrix}g\\\ell_1\end{pmatrix}.
 \quad}
 \label{eq:equalweightlogfusion}
\end{equation}
It has determinant $-1$ and its square is the identity.

The limiting area pairing is
\begin{align}
 \mathcal A(z)
   &=2\pi^2\operatorname{Re}(f\overline g)
     =8\operatorname{Re}(\overline{K(z)}K(1-z)),
       \label{eq:equalweightareaexact}\\
   &=2\pi L|f|^2-2\pi\operatorname{Re}(f\overline{\ell_0}).
       \notag
\end{align}
The Hermitian matrices in the two displayed bases are
\begin{equation}
 H_{fg}=\pi^2\begin{pmatrix}0&1\\1&0\end{pmatrix},\qquad
 H_{\log}=\begin{pmatrix}2\pi L&-\pi\\-\pi&0\end{pmatrix}.
 \label{eq:equalweightHermitian}
\end{equation}
The connection~\eqref{eq:equalweightlogfusion} preserves
$H_{\log}$ exactly.

One can see the cancellation rather than just substitute
into a singular formula. Take
$\mu_1=\mu_3=(1-\delta)/2$ and
$\mu_2=\mu_4=(1+\delta)/2$. Then
\[
 h_0=\frac\pi\delta+\pi L+O(\delta),\qquad
 h_1=-\frac\pi\delta+\pi L+O(\delta),\qquad
 f_1-f_0=\delta\ell_0+O(\delta^2).
\]
The pole terms cancel in~\eqref{eq:screeningareapairing},
leaving~\eqref{eq:equalweightareaexact}. This explains
precisely how the regularized residue sector survives the
coincidence of the ordinary blocks.

For counterclockwise loops about $0$ and $1$, respectively,
the period column $(f,g)^T$ has monodromies
\[
 \begin{pmatrix}1&0\\-2i&1\end{pmatrix},\qquad
 \begin{pmatrix}1&-2i\\0&1\end{pmatrix}.
\]
In the oriented torus-period basis
$\Omega_1=4K(z)$, $\Omega_2=4iK(1-z)$, these become
\begin{equation}
 N_0=\begin{pmatrix}1&0\\2&1\end{pmatrix},\qquad
 N_1=\begin{pmatrix}1&-2\\0&1\end{pmatrix}.
 \label{eq:equalweightmonodromy}
\end{equation}
They give the expected projective $\Gamma(2)$ monodromy.
The channel exchange $z\leftrightarrow1-z$ swaps $f,g$ and
acts on $\tau=\Omega_2/\Omega_1$ by $\tau\mapsto-1/\tau$.
For oriented integral cycles this is a symplectic swap up
to a common phase of the developing differential.

Finally,
\begin{equation}
 \boxed{\quad
 \omega_{\mathrm{DM}}
 =\frac{4\pi^2}{|z(1-z)|^2\mathcal A(z)^2}\dd x\wedge\dd y,
 \qquad
 \mathcal A(z)=8\operatorname{Re}(\overline{K(z)}K(1-z)).
 \quad}
 \label{eq:equalweightdensityexplicit}
\end{equation}
Indeed $W(f,g)=-1/[\pi z(1-z)]$, or equivalently
$\tau'=-i\pi/[4z(1-z)K(z)^2]$.
Near $z=0$,
\[
 \mathcal A(z)=2\pi\log\frac{16}{|z|}
          +O\left(|z|\log\frac1{|z|}\right),\qquad
 \omega_{\mathrm{DM}}\sim
 \frac{\dd x\wedge\dd y}
      {|z|^2\log^2(16/|z|)}.
\]
This gives the integrable cusp density directly in the
cross-ratio coordinate.

\subsection{What the calculation establishes}

The first-moment residue, the screened Virasoro blocks, and the
independently computed Hermitian period pairing agree with all
constants fixed. Under~\eqref{eq:screeningtailinput} and the stated
contour-continuation conditions, the infinite spectral representation
gives the same answer through its two pinches. The finite fusion
transformation on the two-dimensional one-screening residue sector
is the period connection matrix
\eqref{eq:screeningfusionmatrix}, with
\eqref{eq:equalweightlogfusion} at the fully balanced point.
Its independent derivation requires no infinite-contour interchange.
Matrix entries are not
being identified with unspecified pointwise values of a
bare $b$-$6j$ function.

This completes the finite four-point residue calculation.
Section~\ref{sec:normalizedsixj} matches its quantum-group
normalization explicitly, strengthens the Gram-matrix
obstruction, and computes a classical cubic saddle for the
separate compact-tetrahedron degeneration.

